% ============================================================
% Topic 4 -- Số học và Tư duy số -- part d: problems 41--50 (VI)
% Advanced extensions of problems 16--30
% ============================================================

\needspace{9\baselineskip}
\begin{probbox}{41}{\stars{1}}
\textbf{Tổng các số liên tiếp} (mở rộng Bài 16).
\begin{enumerate}[label=\alph*),leftmargin=1.6em,itemsep=1pt]
\item Hãy tìm \emph{tất cả} các cách viết số $100$ thành tổng của ít nhất hai số nguyên dương liên tiếp.
\item Chứng minh rằng số $64$ không thể viết được thành tổng của ít nhất hai số nguyên dương liên tiếp.
\end{enumerate}
\end{probbox}
\needspace{7\baselineskip}
\probsol{
\textbf{Phân tích đề.} Một tổng các số nguyên dương liên tiếp được xác định bởi hai thông tin: \emph{có bao nhiêu số} (gọi là $k$, với $k\ge 2$) và \emph{số đầu tiên} (gọi là $a$, với $a\ge 1$). Ta cần một công thức tính nhanh tổng như vậy, rồi dùng nó để ``thu hẹp'' các khả năng.

\begin{recallbox}\textbf{Tổng các số cách đều} (cách tính của nhà toán học Gauss):\[\text{tổng}=\frac{(\text{số đầu}+\text{số cuối})\times\text{số số hạng}}{2}.\]\end{recallbox}

\begin{proofbox}Viết tổng hai lần, một lần xuôi, một lần ngược, rồi cộng theo cột: ví dụ $S=9+10+\dots+16$ và $S=16+15+\dots+9$. Mỗi cột có tổng $9+16=25$ và có $8$ cột, nên $2S=25\cdot 8$, $S=100$. Lí luận này đúng với mọi dãy cách đều.\end{proofbox}

\step{1}{Lập đẳng thức chính}
Với $k$ số liên tiếp bắt đầu từ $a$: số cuối là $a+k-1$. Gọi $m=(\text{số đầu}+\text{số cuối})=2a+k-1$. Theo công thức Gauss: tổng $=\dfrac{m\cdot k}{2}$. Nếu tổng bằng $100$ thì
\[k\cdot m=200.\]

\step{2}{Hai nhận xét quan trọng về $k$ và $m$}
\begin{itemize}[leftmargin=1.4em,itemsep=1pt]
\item $m-k=2a-1$ là số lẻ, nên \textbf{$k$ và $m$ khác tính chẵn lẻ} (một số chẵn, một số lẻ).
\item $m-k=2a-1\ge 1$, nên \textbf{$m>k$}.
\end{itemize}

\step{3}{Liệt kê các cách phân tích $200=k\cdot m$ với $2\le k<m$}
Vì $k<m$ và $k\cdot m=200$ nên $k\cdot k<200$, tức $k\le 14$. Các ước của $200$ từ $2$ đến $14$ là $2,4,5,8,10$.
\begin{center}
\begin{tabular}{|c|c|c|c|}
\hline
\rowcolor{Tint}$k$ & $m=200:k$ & Khác tính chẵn lẻ? & Kết quả\\
\hline
$2$ & $100$ & không (cùng chẵn) & loại\\
$4$ & $50$ & không & loại\\
$5$ & $40$ & có & $2a+4=40\Rightarrow a=18$\\
$8$ & $25$ & có & $2a+7=25\Rightarrow a=9$\\
$10$ & $20$ & không & loại\\
\hline
\end{tabular}
\end{center}

\step{4}{Viết ra các tổng}
\[100=18+19+20+21+22\qquad\text{và}\qquad 100=9+10+11+12+13+14+15+16.\]
Có đúng \textbf{hai} cách.

\begin{tipbox}Nhớ Bài 16: với $5$ số liên tiếp (số lẻ số hạng), số ở giữa bằng trung bình cộng $=100:5=20$, nên năm số là $18,19,20,21,22$. Ý tưởng ``số ở giữa'' giúp kiểm tra rất nhanh.\end{tipbox}

\step{5}{Câu b): vì sao $64$ không viết được}
Giả sử $64$ là tổng của $k\ge 2$ số nguyên dương liên tiếp. Lập luận như trên: $k\cdot m=2\cdot 64=128$ với $m>k\ge 2$ và $k$, $m$ khác tính chẵn lẻ. Như vậy một trong hai số $k$, $m$ là số lẻ và là ước của $128$. Nhưng $128=2^7$ chỉ có \emph{một} ước lẻ là $1$, trong khi cả $k$ và $m$ đều $\ge 2$. Mâu thuẫn! Vậy $64$ không viết được như yêu cầu.

\begin{proofbox}Cùng lí luận đó chứng tỏ: \emph{mọi lũy thừa của $2$} ($2,4,8,16,32,64,\dots$) đều không viết được thành tổng của ít nhất hai số nguyên dương liên tiếp. Ngược lại, số nào có ước lẻ lớn hơn $1$ thì đều viết được.\end{proofbox}

\begin{warnbox}Đừng quên điều kiện ``nguyên dương'': nếu cho phép số $0$ và số âm thì có thêm những cách ``giả'', ví dụ $(-17)+(-16)+\dots+21+22=100$ (các số từ $-17$ đến $17$ triệt tiêu nhau, còn lại $18+19+20+21+22$). Trong bài này mọi số phải $\ge 1$, nên $a\ge 1$, dẫn tới $m>k$.\end{warnbox}

\verify{$18+19+20+21+22=100$; $9+10+\dots+16=\dfrac{(9+16)\cdot 8}{2}=100$. Thử tay với $64$: $2$ số liên tiếp có tổng lẻ; $3$ số liên tiếp có tổng chia hết cho $3$; \dots\ không cách nào ra $64$.}
\ans{a) Hai cách: $18+19+20+21+22$ và $9+10+\dots+16$. \ b) Xem chứng minh.}
}

\needspace{9\baselineskip}
\begin{probbox}{42}{\stars{1}}
\textbf{Ước chung và số ước} (mở rộng Bài 17).
\begin{enumerate}[label=\alph*),leftmargin=1.6em,itemsep=1pt]
\item Có bao nhiêu số nguyên dương là ước của cả hai số $720$ và $1080$?
\item Tìm số nguyên dương nhỏ nhất có đúng $12$ ước.
\end{enumerate}
\end{probbox}
\needspace{7\baselineskip}
\probsol{
\textbf{Phân tích đề.} Liệt kê ước của $720$ và $1080$ thì rất dài (mỗi số có vài chục ước). Ta dùng hai công cụ mạnh: (1) ước chung của hai số chính là ước của ƯCLN; (2) công thức đếm số ước từ phân tích ra thừa số nguyên tố.

\begin{recallbox}\textbf{Đếm số ước.} Nếu $n=p^{x}\cdot q^{y}\cdot r^{z}$ ($p,q,r$ là các số nguyên tố khác nhau) thì $n$ có $(x+1)(y+1)(z+1)$ ước.\end{recallbox}

\begin{proofbox}Mỗi ước của $n$ có dạng $p^{i}\cdot q^{j}\cdot r^{l}$ với $0\le i\le x$, $0\le j\le y$, $0\le l\le z$. Số mũ $i$ có $x+1$ cách chọn ($0,1,\dots,x$), $j$ có $y+1$ cách, $l$ có $z+1$ cách. Theo quy tắc nhân có $(x+1)(y+1)(z+1)$ ước, và các lựa chọn khác nhau cho các ước khác nhau (vì phân tích ra thừa số nguyên tố là duy nhất).\end{proofbox}

\needspace{6\baselineskip}\par\smallskip\textbf{Câu a).}
\step{1}{Phân tích ra thừa số nguyên tố}
$720=2^4\cdot 3^2\cdot 5$ (vì $720=16\cdot 45=16\cdot 9\cdot 5$); \quad $1080=2^3\cdot 3^3\cdot 5$ (vì $1080=8\cdot 135=8\cdot 27\cdot 5$).

\step{2}{Tìm ƯCLN}
Lấy mỗi thừa số nguyên tố chung với số mũ \emph{nhỏ nhất}: $\text{ƯCLN}(720,1080)=2^3\cdot 3^2\cdot 5=8\cdot 9\cdot 5=360$.
(Kiểm tra bằng cách khác: $1080-720=360$ và $720=2\cdot 360$, nên ƯCLN là $360$.)

\step{3}{Đếm ước của $360$}
$360=2^3\cdot 3^2\cdot 5^1$ có $(3+1)(2+1)(1+1)=4\cdot 3\cdot 2=24$ ước. Vậy có $24$ số là ước của cả $720$ và $1080$.

\needspace{6\baselineskip}\par\smallskip\textbf{Câu b).}
\step{4}{Các ``dạng'' của số có đúng $12$ ước}
Số ước là tích các thừa số (số mũ $+1$). Ta cần viết $12$ thành tích các số $\ge 2$:
\[12=12,\qquad 12=6\cdot 2,\qquad 12=4\cdot 3,\qquad 12=3\cdot 2\cdot 2.\]
Tương ứng với bốn dạng: $p^{11}$, \ $p^{5}q$, \ $p^{3}q^{2}$, \ $p^{2}qr$.

\step{5}{Số nhỏ nhất của mỗi dạng}
Để số nhỏ nhất, dùng các số nguyên tố nhỏ nhất $2,3,5$, và đặt số mũ \emph{lớn} cho số nguyên tố \emph{nhỏ}:
\begin{center}
\begin{tabular}{|c|c|c|}
\hline
\rowcolor{Tint}Dạng & Số nhỏ nhất & Giá trị\\
\hline
$p^{11}$ & $2^{11}$ & $2048$\\
$p^{5}q$ & $2^5\cdot 3$ & $96$\\
$p^{3}q^{2}$ & $2^3\cdot 3^2$ & $72$\\
$p^{2}qr$ & $2^2\cdot 3\cdot 5$ & $\mathbf{60}$\\
\hline
\end{tabular}
\end{center}
Nhỏ nhất là $60$.

\begin{warnbox}Đừng đặt số mũ lớn cho số nguyên tố lớn: ví dụ $2\cdot 3^5=486$ lớn hơn nhiều so với $2^5\cdot 3=96$. Cũng đừng quên dạng có \emph{ba} thừa số nguyên tố ($p^2qr$) -- chính dạng này cho đáp số.\end{warnbox}

\verify{Các ước của $60$: $1,2,3,4,5,6,10,12,15,20,30,60$ -- đúng $12$ ước. Duyệt bằng máy tính các số từ $1$ đến $59$: không số nào có $12$ ước.}
\ans{a) $24$ số. \ b) $60$.}
}

\needspace{9\baselineskip}
\begin{probbox}{43}{\stars{1}}
\textbf{Phân số thay đổi thế nào?} (mở rộng Bài 18 và 19).
\begin{enumerate}[label=\alph*),leftmargin=1.6em,itemsep=1pt]
\item Khi cộng cùng một số tự nhiên $n$ vào cả tử số và mẫu số của phân số $\dfrac58$, ta được một phân số bằng $\dfrac45$. Tìm $n$.
\item Cho phân số $\dfrac ab$ với $0<a<b$. Chứng minh rằng với mọi số nguyên dương $n$ thì $\dfrac{a+n}{b+n}>\dfrac ab$.
\item Có bao nhiêu phân số bằng $\dfrac37$ mà mẫu số là một số có hai chữ số?
\end{enumerate}
\end{probbox}
\needspace{7\baselineskip}
\probsol{
\textbf{Phân tích đề.} Câu a) và b) cùng nói về việc ``cộng thêm cùng một số vào tử và mẫu''. Điều \emph{không thay đổi} khi làm vậy là \emph{hiệu} giữa mẫu và tử -- đây là chìa khóa.

\needspace{6\baselineskip}\par\smallskip\textbf{Câu a).}
\step{1}{Tìm đại lượng không đổi}
Mẫu $-$ tử của $\dfrac58$ là $8-5=3$. Sau khi cộng $n$ vào cả hai: $(8+n)-(5+n)=3$, vẫn bằng $3$.

\step{2}{Lập phương trình}
Sau khi cộng $n$ vào cả tử và mẫu, ta được phân số $\dfrac{5+n}{8+n}$. Theo đề bài:
\[\frac{5+n}{8+n}=\frac45.\]
\begin{recallbox}\textbf{Tính chất của tỉ lệ thức:} nếu $\dfrac ab=\dfrac cd$ (với $b,d\neq 0$) thì $a\cdot d=b\cdot c$ (tích chéo bằng nhau).\end{recallbox}
Áp dụng: $5\cdot(5+n)=4\cdot(8+n)$.

\step{3}{Giải phương trình tìm $n$}
\[25+5n=32+4n\ \Longrightarrow\ 5n-4n=32-25\ \Longrightarrow\ n=7.\]
Khi đó tử mới là $5+7=12$, mẫu mới là $8+7=15$; hiệu vẫn là $15-12=3$ (khớp với Bước 1).

\verify{$\dfrac{5+7}{8+7}=\dfrac{12}{15}=\dfrac{12:3}{15:3}=\dfrac45$. Đúng.}

\needspace{6\baselineskip}\par\smallskip\textbf{Câu b).}
\step{4}{So sánh ``phần còn thiếu để bằng $1$''}
Vì $a<b$ nên cả hai phân số đều nhỏ hơn $1$. Phần còn thiếu để đủ $1$ là
\[1-\frac ab=\frac{b-a}{b},\qquad 1-\frac{a+n}{b+n}=\frac{(b+n)-(a+n)}{b+n}=\frac{b-a}{b+n}.\]
Hai phân số $\dfrac{b-a}{b}$ và $\dfrac{b-a}{b+n}$ có \emph{cùng tử số dương}, còn mẫu $b+n>b$. Phân số có mẫu lớn hơn thì nhỏ hơn, nên $\dfrac{b-a}{b+n}<\dfrac{b-a}{b}$.

\step{5}{Kết luận}
Phân số $\dfrac{a+n}{b+n}$ ``thiếu ít hơn'' để bằng $1$, tức là nó lớn hơn: $\dfrac{a+n}{b+n}>\dfrac ab$.

\begin{recallbox}Hai phân số dương có cùng tử số: phân số nào có mẫu lớn hơn thì nhỏ hơn (cùng một cái bánh chia cho nhiều người hơn thì mỗi phần nhỏ hơn).\end{recallbox}

\begin{proofbox}Cách khác (nhân chéo): với các mẫu dương, $\dfrac{a+n}{b+n}>\dfrac ab$ khi và chỉ khi $(a+n)\cdot b>a\cdot(b+n)$, tức $ab+nb>ab+na$, tức $nb>na$. Điều này đúng vì $n>0$ và $b>a$.\end{proofbox}

\begin{warnbox}Đây chính là lí do ở Bài 18 ta \emph{không được} ``cộng thêm'' vào tử và mẫu: cộng thêm làm phân số thay đổi (nhỏ hơn $1$ thì tăng lên). Chỉ \emph{nhân} (hoặc chia) cả tử và mẫu với cùng một số khác $0$ mới giữ nguyên phân số.\end{warnbox}

\needspace{6\baselineskip}\par\smallskip\textbf{Câu c).}
\step{6}{Dạng chung và điều kiện}
Các phân số bằng $\dfrac37$ có dạng $\dfrac{3k}{7k}$ ($k$ nguyên dương). Mẫu $7k$ có hai chữ số: $10\le 7k\le 99$. Ta có $7\cdot 1=7<10$, $7\cdot 2=14\ge 10$; $7\cdot 14=98\le 99$, $7\cdot 15=105>99$. Vậy $k=2,3,\dots,14$: có $14-2+1=13$ giá trị.
Các phân số là $\dfrac{6}{14},\dfrac{9}{21},\dots,\dfrac{42}{98}$.

\ans{a) $n=7$. \ b) Xem chứng minh. \ c) $13$ phân số.}
}

\needspace{9\baselineskip}
\begin{probbox}{44}{\stars{1}}
\textbf{Cột điện, xe đạp và đoàn tàu} (mở rộng Bài 20).
\begin{enumerate}[label=\alph*),leftmargin=1.6em,itemsep=1pt]
\item Leo đạp xe đều với vận tốc $18$ km/h dọc một con đường có các cột điện cách đều nhau $50$ m. Hỏi từ lúc Leo đi qua cột điện thứ nhất đến lúc Leo đi qua cột điện thứ $25$ mất bao nhiêu phút?
\item Một đoàn tàu chạy đều đi qua một cột điện hết $9$ giây và đi qua một cây cầu dài $350$ m hết $30$ giây (tính từ lúc đầu tàu bắt đầu lên cầu đến lúc đuôi tàu vừa rời khỏi cầu). Tính chiều dài đoàn tàu và vận tốc của tàu theo km/h.
\end{enumerate}
\end{probbox}
\needspace{7\baselineskip}
\probsol{
\textbf{Phân tích đề.} Câu a) có ``bẫy trồng cây'': $25$ cột thì có bao nhiêu khoảng? Câu b) cần hiểu rõ khi ``đi qua'' một vật, đoàn tàu (một vật có chiều dài) phải đi được quãng đường bao nhiêu.

\begin{recallbox}Quãng đường $=$ vận tốc $\times$ thời gian. Đổi $1$ km/h: $1000$ m trong $3600$ giây, nên $1$ km/h $=\dfrac{1000}{3600}=\dfrac{5}{18}$ m/giây; ngược lại $1$ m/giây $=3{,}6$ km/h.\end{recallbox}

\needspace{6\baselineskip}\par\smallskip\textbf{Câu a).}
\step{1}{Đếm số khoảng giữa các cột}
\begin{figC}
\foreach \x in {0,1,2,3}{\draw[Navy,line width=1.5pt] (\x*1.2,0)--(\x*1.2,0.8); \node[below,font=\scriptsize] at (\x*1.2,0) {cột \pgfmathparse{int(\x+1)}\pgfmathresult};}
\foreach \x in {0,1,2}{\draw[<->,tmOrange] (\x*1.2+0.05,0.4)--(\x*1.2+1.15,0.4); \node[above,font=\scriptsize,text=tmOrange!80!black] at (\x*1.2+0.6,0.4) {$50$ m};}
\node at (4.35,0.4) {$\cdots$};
\node[anchor=west,font=\footnotesize,text=Ink] at (5,0.4) {$4$ cột $\to$ $3$ khoảng};
\end{figC}
Giữa cột thứ nhất và cột thứ $25$ có $25-1=24$ khoảng, tổng chiều dài $24\cdot 50=1200$ m.

\step{2}{Đổi vận tốc và tính thời gian}
$18$ km/h $=18\cdot\dfrac{5}{18}=5$ m/giây. Thời gian: $1200:5=240$ giây $=240:60=4$ phút.

\begin{warnbox}Nếu tính $25$ khoảng ($25\cdot 50=1250$ m) sẽ được $250$ giây -- sai. Số khoảng luôn ít hơn số cột $1$ đơn vị (khi tính từ cột đầu đến cột cuối).\end{warnbox}

\needspace{6\baselineskip}\par\smallskip\textbf{Câu b).}
\step{3}{Hiểu ``đi qua cột điện'' và ``đi qua cây cầu''}
\begin{itemize}[leftmargin=1.4em,itemsep=1pt]
\item Đi qua cột điện: từ lúc đầu tàu ngang cột đến lúc đuôi tàu ngang cột, tàu đi được đúng \emph{một chiều dài tàu} $L$.
\item Đi qua cầu: từ lúc đầu tàu ở đầu cầu đến lúc đuôi tàu rời cuối cầu, tàu đi được \emph{chiều dài cầu cộng chiều dài tàu}: $350+L$.
\end{itemize}
\begin{figC}
% scale: 1 unit = 70 m (bridge 350 m = 5 units, train 150 m = 15/7 units)
\pgfmathsetmacro{\rvdbL}{15/7}
\fill[Slate!25] ({\rvdbL},-0.15) rectangle ({\rvdbL+5},0); \node[below,font=\scriptsize] at ({\rvdbL+2.5},-0.15) {cầu $350$ m};
\draw[tmBlue,fill=tmBlue!20,rounded corners=2pt] (0,0.05) rectangle ({\rvdbL},0.45); \node[font=\scriptsize] at ({\rvdbL/2},0.25) {tàu (lúc đầu)};
\draw[tmOrange,fill=tmOrange!20,rounded corners=2pt] ({\rvdbL+5},0.05) rectangle ({2*\rvdbL+5},0.45); \node[font=\scriptsize] at ({1.5*\rvdbL+5},0.25) {tàu (lúc cuối)};
\draw[<->,Slate] (0,-0.3)--({\rvdbL},-0.3) node[midway,below,font=\scriptsize] {dài $L$};
\draw[dashed,Slate] ({\rvdbL},0.05)--({\rvdbL},0.8); \draw[dashed,Slate] ({2*\rvdbL+5},0.05)--({2*\rvdbL+5},0.8);
\draw[<->,Navy] ({\rvdbL},0.8)--({2*\rvdbL+5},0.8) node[midway,above,font=\scriptsize] {đầu tàu đi được $350+L$};
\end{figC}

\step{4}{So sánh hai lần}
Qua cầu mất $30$ giây, qua cột mất $9$ giây. Thời gian chênh lệch $30-9=21$ giây chính là thời gian tàu đi thêm $350$ m (chiều dài cầu). Vận tốc tàu:
\[v=350:21=\frac{350}{21}=\frac{50}{3}\ \text{m/giây}.\]

\step{5}{Tính chiều dài tàu và đổi vận tốc}
Chiều dài tàu $L=v\cdot 9=\dfrac{50}{3}\cdot 9=150$ m. Vận tốc: $\dfrac{50}{3}\cdot 3{,}6=\dfrac{50\cdot 3{,}6}{3}=60$ km/h.

\begin{warnbox}Lỗi hay gặp: cho rằng qua cầu tàu chỉ đi $350$ m (quên cộng chiều dài tàu), khi đó vận tốc là $350:30\approx 11{,}7$ m/giây và kết quả sai. Hãy theo dõi \emph{một điểm cố định} trên tàu (ví dụ đầu tàu) để biết tàu đã đi bao xa.\end{warnbox}

\verify{Với $v=\dfrac{50}{3}$ m/giây và $L=150$ m: qua cột mất $150:\dfrac{50}{3}=9$ giây; qua cầu mất $(350+150):\dfrac{50}{3}=500\cdot\dfrac{3}{50}=30$ giây. Đúng cả hai dữ kiện.}
\ans{a) $4$ phút. \ b) Tàu dài $150$ m, vận tốc $60$ km/h.}
}

\needspace{9\baselineskip}
\begin{probbox}{45}{\stars{2}}
\textbf{Tích các chữ số} (mở rộng Bài 21).
\begin{enumerate}[label=\alph*),leftmargin=1.6em,itemsep=1pt]
\item Có bao nhiêu số có bốn chữ số mà tích các chữ số bằng $36$?
\item Tìm số nguyên dương nhỏ nhất có tích các chữ số bằng $360$.
\item Tìm số nguyên dương lớn nhất có các chữ số \emph{đôi một khác nhau} và tích các chữ số bằng $360$.
\end{enumerate}
\end{probbox}
\needspace{7\baselineskip}
\probsol{
\textbf{Phân tích đề.} Cả ba câu đều bắt đầu bằng việc \emph{phân tích số ra thừa số nguyên tố}, rồi ``chia'' các thừa số đó vào các chữ số (mỗi chữ số từ $1$ đến $9$). Chú ý: các chữ số đều khác $0$ (nếu không tích bằng $0$).

\begin{recallbox}$36=2^2\cdot 3^2$; \ $360=2^3\cdot 3^2\cdot 5$. Các chữ số và phân tích của chúng: $1$; $2$; $3$; $4=2^2$; $5$; $6=2\cdot 3$; $7$; $8=2^3$; $9=3^2$.\end{recallbox}

\needspace{6\baselineskip}\par\smallskip\textbf{Câu a).}
\step{1}{Tìm các bộ bốn chữ số (không kể thứ tự) có tích $36$}
Làm như Bài 21, xếp $a\le b\le c\le d$:
\begin{itemize}[leftmargin=1.4em,itemsep=1pt]
\item Có hai chữ số $1$: còn lại tích $36$ với hai chữ số: $\{1,1,4,9\}$, $\{1,1,6,6\}$.
\item Có đúng một chữ số $1$: ba chữ số còn lại ($\ge 2$) có tích $36$: $\{1,2,2,9\}$, $\{1,2,3,6\}$, $\{1,3,3,4\}$.
\item Không có chữ số $1$: bốn chữ số $\ge 2$ có tích $36$: chỉ có $\{2,2,3,3\}$ (vì $36=2\cdot2\cdot3\cdot3$, không thể tách thêm).
\item Có ba chữ số $1$: chữ số còn lại phải là $36$ -- không được.
\end{itemize}

\begin{recallbox}\textbf{Đếm cách xếp $4$ chữ số:} cả bốn khác nhau: $4\cdot 3\cdot 2\cdot 1=24$ cách. Có đúng hai chữ số giống nhau (hai cái còn lại khác nhau và khác chúng): chọn $2$ vị trí cho cặp giống nhau trong $4$ vị trí ($6$ cách), rồi xếp hai chữ số còn lại vào $2$ chỗ trống ($2$ cách): $12$ cách. Hai cặp giống nhau (như $1,1,6,6$): chỉ cần chọn $2$ vị trí cho một cặp: $6$ cách.\end{recallbox}

\step{2}{Đếm}
\begin{center}
\begin{tabular}{|c|c|c|}
\hline
\rowcolor{Tint}Bộ chữ số & Loại & Số cách xếp\\
\hline
$\{1,1,4,9\}$ & một cặp giống & $12$\\
$\{1,1,6,6\}$ & hai cặp giống & $6$\\
$\{1,2,2,9\}$ & một cặp giống & $12$\\
$\{1,2,3,6\}$ & khác nhau & $24$\\
$\{1,3,3,4\}$ & một cặp giống & $12$\\
$\{2,2,3,3\}$ & hai cặp giống & $6$\\
\hline
\multicolumn{2}{|r|}{\textbf{Tổng}} & $\mathbf{72}$\\
\hline
\end{tabular}
\end{center}

\needspace{6\baselineskip}\par\smallskip\textbf{Câu b).}
\step{3}{Cần ít chữ số nhất}
Số có ít chữ số hơn thì nhỏ hơn. Một hoặc hai chữ số có tích lớn nhất $9\cdot 9=81<360$, nên cần ít nhất $3$ chữ số.

\step{4}{Tìm bộ ba chữ số có tích $360$}
$360$ có thừa số nguyên tố $5$, mà chữ số duy nhất chứa thừa số $5$ là chính chữ số $5$. Vậy một chữ số là $5$, hai chữ số còn lại có tích $360:5=72=8\cdot 9$ (cách duy nhất với các chữ số $\le 9$). Bộ chữ số: $\{5,8,9\}$.

\step{5}{Xếp để được số nhỏ nhất}
Xếp các chữ số từ nhỏ đến lớn: $589$.

\needspace{6\baselineskip}\par\smallskip\textbf{Câu c).}
\step{6}{Muốn số lớn nhất: càng nhiều chữ số càng tốt}
Các chữ số khác nhau, đều khác $0$, và một chữ số phải là $5$. Sáu chữ số khác nhau thì tích ít nhất là $1\cdot2\cdot3\cdot4\cdot5\cdot6=720>360$, nên có nhiều nhất $5$ chữ số.

\step{7}{Tìm các bộ $5$ chữ số khác nhau có tích $360$}
Một chữ số là $5$; bốn chữ số còn lại khác nhau, khác $5$, có tích $72=2^3\cdot 3^2$. Bốn chữ số khác nhau có tích $72$ mà nhỏ như vậy thì phải có chữ số $1$ (vì $2\cdot3\cdot4\cdot6=144>72$); ba chữ số còn lại khác nhau có tích $72$: $72=2\cdot4\cdot9=3\cdot4\cdot6$ (còn $2\cdot 6\cdot 6$, $3\cdot 3\cdot 8$ có chữ số lặp; $1\cdot 8\cdot 9$ lặp chữ số $1$). Được hai bộ: $\{1,2,4,5,9\}$ và $\{1,3,4,5,6\}$.

\step{8}{Xếp để được số lớn nhất}
Xếp từ lớn đến nhỏ: $95421$ và $65431$. Số lớn hơn là $95421$.

\begin{warnbox}Ở câu c), đừng chọn ngay số có chữ số đầu lớn như $985$ (tích $360$): số có $5$ chữ số luôn lớn hơn số có $3$ chữ số. Hãy ưu tiên \emph{số chữ số} trước, rồi mới đến chữ số đầu.\end{warnbox}

\verify{$5\cdot 8\cdot 9=360$; $9\cdot5\cdot4\cdot2\cdot1=360$. Máy tính duyệt mọi số có tới $6$ chữ số xác nhận: có $72$ số bốn chữ số tích $36$; số nhỏ nhất tích $360$ là $589$; số lớn nhất có chữ số khác nhau, tích $360$ là $95\,421$.}
\ans{a) $72$ số. \ b) $589$. \ c) $95\,421$.}
}

\needspace{22\baselineskip}
\begin{probbox}{46}{\stars{1}}
\textbf{Bảng cách đều} (mở rộng Bài 22). Trong một bảng ô vuông, các số trên mỗi hàng và trên mỗi cột đều tạo thành dãy cách đều.
\begin{enumerate}[label=\alph*),leftmargin=1.6em,itemsep=1pt]
\item Bảng $4\times 4$ có bốn số ở bốn góc là: góc trên trái $1$, góc trên phải $10$, góc dưới trái $19$, góc dưới phải $37$. Hãy điền đầy đủ bảng.
\item Tính tổng của tất cả $16$ số trong bảng ở câu a). So sánh với tổng bốn số ở góc, em thấy gì? Giải thích.
\item Chứng minh rằng trong một bảng $3\times 3$ như vậy, số ở ô chính giữa luôn bằng trung bình cộng của bốn số ở bốn góc.
\end{enumerate}
\begin{figC}
\foreach \rr in {0,1,2,3}{\foreach \cc in {0,1,2,3}{\draw[Navy,line width=0.9pt,fill=white] (\cc*1.1,-\rr*0.8) rectangle (\cc*1.1+1.1,-\rr*0.8-0.8);}}
\node at (0.55,-0.4) {$1$}; \node at (3.85,-0.4) {$10$};
\node at (0.55,-2.8) {$19$}; \node at (3.85,-2.8) {$37$};
\end{figC}
\end{probbox}
\needspace{7\baselineskip}
\probsol{
\textbf{Phân tích đề.} Bảng chỉ cho bốn số ở góc. Cột đầu và cột cuối (cũng như hàng đầu, hàng cuối) mỗi cái đã biết số đầu và số cuối, nên tìm được ngay. Sau đó dùng các hàng hoặc các cột ở giữa.

\begin{recallbox}Dãy cách đều $4$ số hạng: khoảng cách $d=(\text{số cuối}-\text{số đầu}):3$. Tổng $4$ số hạng $=2\times(\text{số đầu}+\text{số cuối})$ (ghép số $1$ với số $4$, số $2$ với số $3$: hai cặp có tổng bằng nhau). Dãy $3$ số hạng: số giữa $=$ trung bình cộng hai số ngoài.\end{recallbox}

\needspace{6\baselineskip}\par\smallskip\textbf{Câu a).}
\step{1}{Hàng đầu và hàng cuối}
Hàng $1$: từ $1$ đến $10$, $d=(10-1):3=3$: \ $1,4,7,10$.\\
Hàng $4$: từ $19$ đến $37$, $d=(37-19):3=6$: \ $19,25,31,37$.

\step{2}{Bốn cột (biết số đầu và số cuối của mỗi cột)}
Cột $1$: $1\to 19$, $d=6$: $1,7,13,19$. \quad Cột $2$: $4\to 25$, $d=7$: $4,11,18,25$.\\
Cột $3$: $7\to 31$, $d=8$: $7,15,23,31$. \quad Cột $4$: $10\to 37$, $d=9$: $10,19,28,37$.

\step{3}{Kiểm tra các hàng ở giữa}
Hàng $2$: $7,11,15,19$ (cách đều $4$); hàng $3$: $13,18,23,28$ (cách đều $5$). Bảng thỏa mãn mọi điều kiện:
\begin{figC}
\foreach \rr in {0,1,2,3}{\foreach \cc in {0,1,2,3}{\draw[Navy,line width=0.9pt,fill=white] (\cc*1.1,-\rr*0.8) rectangle (\cc*1.1+1.1,-\rr*0.8-0.8);}}
\foreach \cc/\val in {0/1,1/4,2/7,3/10}{\node at (\cc*1.1+0.55,-0.4) {$\val$};}
\foreach \cc/\val in {0/7,1/11,2/15,3/19}{\node at (\cc*1.1+0.55,-1.2) {$\val$};}
\foreach \cc/\val in {0/13,1/18,2/23,3/28}{\node at (\cc*1.1+0.55,-2.0) {$\val$};}
\foreach \cc/\val in {0/19,1/25,2/31,3/37}{\node at (\cc*1.1+0.55,-2.8) {$\val$};}
\foreach \rr/\val in {0/22,1/52,2/82,3/112}{\node[anchor=west,font=\footnotesize,text=Slate] at (4.55,-\rr*0.8-0.4) {tổng hàng: $\val$};}
\end{figC}

\needspace{6\baselineskip}\par\smallskip\textbf{Câu b).}
\step{4}{Tính tổng}
Tổng các hàng: $22+52+82+112=268$. Tổng bốn góc: $1+10+19+37=67$. Ta thấy $268=4\cdot 67$: \textbf{tổng cả bảng bằng $4$ lần tổng bốn góc.}

\step{5}{Giải thích (đúng với mọi bảng $4\times4$ như vậy)}
Mỗi hàng là dãy cách đều $4$ số, nên tổng mỗi hàng $=2\times$(số đầu hàng $+$ số cuối hàng). Cộng $4$ hàng:
\[\text{Tổng bảng}=2\times(\text{tổng cột }1+\text{tổng cột }4).\]
Cột $1$ và cột $4$ cũng là dãy cách đều $4$ số: tổng cột $1=2\times(1+19)$, tổng cột $4=2\times(10+37)$. Vậy
\[\text{Tổng bảng}=2\times\big(2\cdot(1+19)+2\cdot(10+37)\big)=4\times(1+10+19+37)=4\times 67=268.\]

\needspace{6\baselineskip}\par\smallskip\textbf{Câu c).}
\step{6}{Đặt tên các ô của bảng $3\times 3$}
\begin{figC}
\foreach \rr in {0,1,2}{\foreach \cc in {0,1,2}{\draw[Navy,line width=0.9pt,fill=white] (\cc*1.1,-\rr*0.8) rectangle (\cc*1.1+1.1,-\rr*0.8-0.8);}}
\node at (0.55,-0.4) {$a$}; \node at (1.65,-0.4) {$\cdot$}; \node at (2.75,-0.4) {$b$};
\node at (0.55,-1.2) {$u$}; \node[text=AnsRed] at (1.65,-1.2) {$x$}; \node at (2.75,-1.2) {$v$};
\node at (0.55,-2.0) {$c$}; \node at (1.65,-2.0) {$\cdot$}; \node at (2.75,-2.0) {$d$};
\end{figC}
Bốn góc là $a,b,c,d$; ô giữa là $x$; $u$ và $v$ là ô giữa của cột trái và cột phải.

\step{7}{Dùng tính chất ``số giữa $=$ trung bình cộng hai số ngoài''}
Cột trái $a,u,c$ cách đều nên $u=\dfrac{a+c}{2}$. Cột phải: $v=\dfrac{b+d}{2}$. Hàng giữa $u,x,v$ cách đều nên
\[x=\frac{u+v}{2}=\frac{\frac{a+c}{2}+\frac{b+d}{2}}{2}=\frac{a+b+c+d}{4}.\]
Đó chính là trung bình cộng của bốn số ở góc.

\begin{proofbox}Vì sao trong dãy cách đều $p,q,r$ ta có $q=\dfrac{p+r}{2}$? Vì $q-p=r-q$ (cùng bằng khoảng cách $d$), nên $2q=p+r$.\end{proofbox}

\verify{Áp dụng câu c) cho phần $3\times 3$ ở góc trên trái của bảng câu a): bốn góc $1,7,13,23$, trung bình cộng $\dfrac{1+7+13+23}{4}=11$, đúng bằng ô giữa $11$.}
\ans{a) Xem bảng. \ b) Tổng $268=4\times 67$ (gấp $4$ lần tổng bốn góc). \ c) Xem chứng minh.}
}

\needspace{9\baselineskip}
\begin{probbox}{47}{\stars{2}}
\textbf{Số ``xoay vòng''} (mở rộng Bài 23). Trong các câu dưới đây, mỗi chữ cái biểu thị một chữ số, và $A\neq 0$.
\begin{enumerate}[label=\alph*),leftmargin=1.6em,itemsep=1pt]
\item Tìm số có sáu chữ số $\nd{2ABCDE}$ sao cho $\nd{2ABCDE}\times 3=\nd{ABCDE2}$.
\item Tìm tất cả các chữ số $d$ (từ $1$ đến $9$) sao cho có số có sáu chữ số $\nd{dABCDE}$ thỏa mãn $\nd{dABCDE}\times 3=\nd{ABCDEd}$.
\item Chứng minh rằng không có số có sáu chữ số $\nd{1ABCDE}$ nào thỏa mãn $\nd{1ABCDE}\times 2=\nd{ABCDE1}$.
\end{enumerate}
\end{probbox}
\needspace{7\baselineskip}
\probsol{
\textbf{Phân tích đề.} Cả ba câu đều là ``chuyển chữ số đầu xuống cuối''. Ta dùng lại Cách 1 của Bài 23: coi khối $\nd{ABCDE}$ là một số $N$ rồi lập đẳng thức. Lần này ta làm một lần cho mọi chữ số đầu $d$.

\begin{recallbox}Với $N=\nd{ABCDE}$: \ $\nd{dABCDE}=d\cdot 100\,000+N$ \ và \ $\nd{ABCDEd}=10\cdot N+d$. Vì $A\neq 0$ nên $N$ là số có đúng năm chữ số: $10\,000\le N\le 99\,999$.\end{recallbox}

\step{1}{Lập đẳng thức tổng quát cho câu a) và b)}
\[3\cdot(d\cdot 100\,000+N)=10N+d\ \Longrightarrow\ 300\,000\,d+3N=10N+d\ \Longrightarrow\ 299\,999\,d=7N.\]
Vì $299\,999=7\cdot 42\,857$ nên chia hai vế cho $7$:
\[N=42\,857\cdot d.\]

\step{2}{Câu a): $d=2$}
$N=42\,857\cdot 2=85\,714$. Số cần tìm là $285\,714$.

\step{3}{Câu b): điều kiện để $N$ có năm chữ số}
$N=42\,857\,d\le 99\,999$. Với $d=1$: $N=42\,857$ (năm chữ số, được). Với $d=2$: $N=85\,714$ (được). Với $d=3$: $N=128\,571$ có sáu chữ số -- quá lớn, không thể là khối $\nd{ABCDE}$. Với $d\ge 3$ thì $N$ còn lớn hơn nữa. Vậy chỉ có $d=1$ và $d=2$, ứng với hai số $142\,857$ và $285\,714$.

\begin{tipbox}Hai số này có cùng các chữ số $1,4,2,8,5,7$ xếp vòng tròn. Đó không phải ngẫu nhiên: $\dfrac17=0{,}(142857)$ và $\dfrac27=0{,}(285714)$. Nhân với $3$ ứng với việc ``xoay'' vòng tròn chữ số.\end{tipbox}

\step{4}{Câu c): lập đẳng thức}
Với $N=\nd{ABCDE}$:
\[2\cdot(100\,000+N)=10N+1\ \Longrightarrow\ 200\,000+2N=10N+1\ \Longrightarrow\ 199\,999=8N.\]

\step{5}{Câu c): tìm mâu thuẫn chẵn -- lẻ}
Vế phải $8N$ là số chẵn (vì $8$ chẵn). Vế trái $199\,999$ là số lẻ. Một số lẻ không thể bằng một số chẵn. Vậy không có số nào thỏa mãn.

\begin{proofbox}Có thể thấy ngay mà không cần lập đẳng thức: $\nd{1ABCDE}\times 2$ là số chẵn, nên chữ số tận cùng của nó phải chẵn; nhưng $\nd{ABCDE1}$ tận cùng bằng $1$ (lẻ). Mâu thuẫn.\end{proofbox}

\begin{warnbox}Ở câu b), nhiều bạn dừng lại khi tìm được $d=1$ (đáp án của Bài 23) và quên $d=2$; hoặc cho $d=3,4,\dots$ cũng được mà không kiểm tra $N$ có quá năm chữ số. Luôn kiểm tra điều kiện ``$N$ có đúng năm chữ số''.\end{warnbox}

\verify{$142\,857\times 3=428\,571$; \ $285\,714\times 3=857\,142$. Cả hai đều đúng là ``chữ số đầu chuyển xuống cuối''.}
\ans{a) $285\,714$. \ b) $d=1$ hoặc $d=2$ (các số $142\,857$ và $285\,714$). \ c) Xem chứng minh.}
}

\needspace{9\baselineskip}
\begin{probbox}{48}{\stars{2}}
\textbf{Chú ếch nhảy trên trục số} (mở rộng Bài 24). Một chú ếch đứng ở điểm $0$ trên trục số. Ở lần nhảy thứ $k$ ($k=1,2,3,\dots$), ếch nhảy đúng $k$ đơn vị, sang trái hoặc sang phải tùy ý.
\begin{enumerate}[label=\alph*),leftmargin=1.6em,itemsep=1pt]
\item Sau đúng $10$ lần nhảy, ếch có thể trở về điểm $0$ không? Vì sao?
\item Sau đúng $10$ lần nhảy, ếch có thể đứng ở những điểm nào? Có tất cả bao nhiêu điểm như vậy?
\item Ếch cần ít nhất bao nhiêu lần nhảy để đến được điểm $20$?
\end{enumerate}
\end{probbox}
\needspace{7\baselineskip}
\probsol{
\textbf{Phân tích đề.} Vị trí của ếch sau $n$ lần nhảy là $\pm1\pm2\pm3\cdots\pm n$ (dấu $+$ khi nhảy sang phải, dấu $-$ khi nhảy sang trái). Ý tưởng giống Bài 24: bắt đầu từ ``tất cả đều cộng'' rồi xem đổi dấu làm thay đổi kết quả thế nào.

\begin{recallbox}$1+2+\dots+n=\dfrac{n(n+1)}{2}$ (công thức Gauss). Ví dụ $1+2+\dots+10=\dfrac{10\cdot 11}{2}=55$.\end{recallbox}

\step{1}{Quy tắc ``đổi dấu''}
Nếu ếch nhảy hết sang phải, sau $10$ lần nó ở điểm $55$. Đổi một lần nhảy $k$ từ phải sang trái làm vị trí giảm $2k$ (mất $k$ và còn lùi thêm $k$). Vậy nếu tổng các bước nhảy sang trái là $s$ thì vị trí cuối là
\[55-2s.\]

\needspace{6\baselineskip}\par\smallskip\textbf{Câu a).}
\step{2}{Tính chẵn lẻ}
$55-2s$ luôn là số \emph{lẻ} (số lẻ trừ số chẵn). Điểm $0$ là số chẵn, nên ếch \textbf{không thể} về điểm $0$ sau $10$ lần nhảy.

\needspace{6\baselineskip}\par\smallskip\textbf{Câu b).}
\step{3}{$s$ có thể nhận những giá trị nào?}
$s$ là tổng của một nhóm (có thể rỗng) các số khác nhau trong $\{1,2,\dots,10\}$, nên $0\le s\le 55$. Ta chứng tỏ \emph{mọi} số nguyên $s$ từ $0$ đến $55$ đều đạt được.

\begin{proofbox}Lấy dần các số lớn nhất: $10, 9, 8, \dots$ Các tổng $0$, $10$, $10+9=19$, $19+8=27$, $27+7=34$, $34+6=40$, $40+5=45$, $45+4=49$, $49+3=52$, $52+2=54$, $54+1=55$ chia đoạn $[0;55]$ thành các khúc. Một số $s$ nằm giữa hai mốc liên tiếp, ví dụ $27<s\le 34$, thì viết $s=10+9+8+(s-27)$ với $s-27$ từ $1$ đến $7$ -- là một số chưa dùng. Tương tự cho mọi khúc khác (với $0<s\le 10$ chỉ cần lấy số $s$). Vậy mọi $s$ từ $0$ đến $55$ đều đạt được.\end{proofbox}

\step{4}{Các vị trí cuối cùng}
Với $s=0,1,2,\dots,55$, vị trí $55-2s$ lần lượt là $55,53,51,\dots,-53,-55$: tức là \textbf{mọi số lẻ từ $-55$ đến $55$}. Số các điểm: từ $-55$ đến $55$ có $28$ số lẻ dương ($1,3,\dots,55$) và $28$ số lẻ âm, tổng cộng $56$ điểm (cũng bằng số giá trị của $s$: $56$).

\begin{figC}
\begin{scope}[x=0.5cm]
\draw[->,thick,Navy] (-12.4,0)--(12.4,0);
\foreach \x in {-11,...,11}{\draw[Navy] (\x,0.07)--(\x,-0.07);}
\node[below,font=\scriptsize] at (0,-0.1) {$0$};
\foreach \x in {-11,-9,...,11}{\fill[KidOrange] (\x,0) circle (2.3pt); \node[below,font=\scriptsize] at (\x,-0.1) {$\x$};}
\node[above,font=\scriptsize,text=Slate] at (0,0.15) {các vị trí lẻ từ $-11$ đến $11$; tiếp tục như vậy tới $-55$ và $55$};
\end{scope}
\end{figC}

\needspace{6\baselineskip}\par\smallskip\textbf{Câu c).}
\step{5}{Hai điều kiện cần}
Sau $n$ lần nhảy, gọi $T=1+2+\dots+n$. Vị trí là $T-2s$, nên: (1) $T\ge 20$ (không thể đi xa hơn $T$); (2) $T-20$ phải chẵn, tức $T$ và $20$ cùng chẵn.

\step{6}{Thử lần lượt $n$}
\begin{center}
\begin{tabular}{|c|c|c|c|c|c|c|c|}
\hline
\rowcolor{Tint}$n$ & $1$ & $2$ & $3$ & $4$ & $5$ & $6$ & $7$\\
\hline
$T$ & $1$ & $3$ & $6$ & $10$ & $15$ & $21$ & $28$\\
\hline
Đạt $20$? & \multicolumn{5}{c|}{không ($T<20$)} & không ($T$ lẻ) & \textbf{có thể}\\
\hline
\end{tabular}
\end{center}

\step{7}{Chỉ ra cách nhảy với $n=7$}
Cần $28-2s=20$, tức $s=4$: chỉ cần nhảy sang trái ở lần thứ $4$:
\[1+2+3-4+5+6+7=20.\]
Vậy ít nhất là $7$ lần nhảy.

\begin{warnbox}Không đủ chỉ kiểm tra ``$T\ge 20$'': với $n=6$ thì $T=21\ge 20$ nhưng vị trí luôn lẻ, không thể là $20$. Đây chính là ý tưởng chẵn -- lẻ của Bài 24.\end{warnbox}

\verify{$1+2+3-4+5+6+7=20$. Với $n=6$: vị trí luôn là $21-2s$, là số lẻ, khác $20$. Máy tính duyệt mọi cách nhảy cũng xác nhận: sau $10$ lần có đúng $56$ vị trí, không có $0$; số lần ít nhất để tới $20$ là $7$.}
\ans{a) Không (vị trí luôn lẻ). \ b) Mọi số lẻ từ $-55$ đến $55$: $56$ điểm. \ c) $7$ lần.}
}

\needspace{9\baselineskip}
\begin{probbox}{49}{\stars{1}}
\textbf{Số chứa khối $2025$} (mở rộng Bài 28). Ta xét các số có sáu chữ số chứa bốn chữ số $2,0,2,5$ \emph{đứng liền nhau theo đúng thứ tự này} (ví dụ: $720\,251$, $202\,513$, $132\,025$). Có bao nhiêu số như vậy chia hết cho $45$?
\end{probbox}
\needspace{7\baselineskip}
\probsol{
\textbf{Phân tích đề.} Số có dạng khối $2025$ cộng thêm hai chữ số $a,b$ ở các vị trí khác nhau. ``Chia hết cho $45$'' cần được tách thành hai điều kiện dễ kiểm tra hơn.

\begin{recallbox}$45=5\cdot 9$, và $5$, $9$ không có ước chung nào khác $1$. Vì vậy một số chia hết cho $45$ khi và chỉ khi nó chia hết cho cả $5$ và $9$.
\textbf{Chia hết cho $5$:} chữ số tận cùng là $0$ hoặc $5$.
\textbf{Chia hết cho $9$:} tổng các chữ số chia hết cho $9$.\end{recallbox}

\step{1}{Điều kiện chia hết cho $9$}
Tổng chữ số của khối: $2+0+2+5=9$. Tổng chữ số của cả số là $9+a+b$, chia hết cho $9$ khi $a+b$ chia hết cho $9$, tức $a+b\in\{0;\ 9;\ 18\}$.

\step{2}{Ba vị trí của khối và đếm từng trường hợp}
\begin{itemize}[leftmargin=1.4em,itemsep=2pt]
\item \textbf{Dạng $\nd{2025ab}$} (chữ số tận cùng là $b$, nên $b\in\{0;5\}$; $a$ tùy ý):
  $b=0$: $a\in\{0;9\}$ $\to$ $202\,500$, $202\,590$. \ $b=5$: $a+5\in\{9;18\}$ $\to$ $a=4$ $\to$ $202\,545$. \ \textbf{Có $3$ số.}
\item \textbf{Dạng $\nd{a2025b}$} ($a\neq 0$, $b\in\{0;5\}$):
  $b=0$: $a=9$ $\to$ $920\,250$. \ $b=5$: $a=4$ $\to$ $420\,255$. \ \textbf{Có $2$ số.}
\item \textbf{Dạng $\nd{ab2025}$} (tận cùng là $5$ -- luôn chia hết cho $5$; $a\neq 0$):
  $a+b=9$: $a=1,2,\dots,9$ (với $b=9-a$) $\to$ $9$ số: $182\,025$, $272\,025$, \dots, $902\,025$. \ $a+b=18$: $a=b=9$ $\to$ $992\,025$. \ ($a+b=0$ không được vì $a\neq0$.) \ \textbf{Có $10$ số.}
\end{itemize}

\step{3}{Kiểm tra đếm trùng}
Một số bị đếm hai lần nếu khối $2025$ xuất hiện hai lần trong sáu chữ số, tức hai khối chồng lên nhau. Dịch khối đi $1$ vị trí: cần ``$025$'' trùng ``$202$'' -- sai; dịch $2$ vị trí: cần ``$25$'' trùng ``$20$'' -- sai. Vậy không có số nào bị đếm trùng.

\step{4}{Cộng lại}
$3+2+10=15$ số.

\begin{warnbox}(1) Đừng quên trường hợp $a+b=0$ ở dạng $\nd{2025ab}$ (số $202\,500$). (2) Ở dạng $\nd{a2025b}$ và $\nd{ab2025}$, chữ số đầu $a$ phải khác $0$. (3) Phải kiểm tra \emph{cả hai} điều kiện chia hết cho $5$ và cho $9$, không chỉ một.\end{warnbox}

\verify{Ví dụ $992\,025$: tận cùng $5$, tổng chữ số $9+9+2+0+2+5=27$ chia hết cho $9$, nên chia hết cho $45$ ($992\,025=45\cdot 22\,045$). Máy tính duyệt mọi số có sáu chữ số cũng được đúng $15$ số.}
\ans{$15$ số.}
}

\needspace{9\baselineskip}
\begin{probbox}{50}{\stars{2}}
\textbf{Trung bình cộng và trung vị} (mở rộng Bài 30). Một danh sách gồm năm số nguyên dương (có thể có số trùng nhau) có trung bình cộng bằng $12$ và trung vị bằng $10$.
\begin{enumerate}[label=\alph*),leftmargin=1.6em,itemsep=1pt]
\item Số lớn nhất trong danh sách có thể lớn nhất là bao nhiêu?
\item Số lớn nhất trong danh sách có thể nhỏ nhất là bao nhiêu?
\item Trả lời lại hai câu hỏi trên nếu năm số đôi một khác nhau.
\end{enumerate}
\end{probbox}
\needspace{7\baselineskip}
\probsol{
\textbf{Phân tích đề.} Xếp năm số từ nhỏ đến lớn: $a\le b\le c\le d\le e$. Trung vị là số giữa nên $c=10$. Trung bình cộng $12$ nên tổng là $5\cdot 12=60$. Ta cần làm cho $e$ lớn nhất hoặc nhỏ nhất. Đây là bài toán ``tối ưu'': phải vừa \emph{chứng minh} không thể vượt qua một giới hạn, vừa \emph{chỉ ra ví dụ} đạt đúng giới hạn đó.

\begin{recallbox}Trung bình cộng $=$ tổng $:$ số các số, nên tổng $=12\cdot 5=60$. Trung vị của năm số (đã xếp tăng dần) là số thứ ba.\end{recallbox}

\step{1}{Đẳng thức chính}
\[a+b+10+d+e=60\quad\Longrightarrow\quad a+b+d+e=50,\qquad a\le b\le 10\le d\le e.\]

\needspace{6\baselineskip}\par\smallskip\textbf{Câu a): làm $e$ lớn nhất.}
\step{2}{Giới hạn trên}
Muốn $e$ lớn thì các số khác phải nhỏ nhất có thể: $a\ge 1$, $b\ge 1$, $d\ge 10$. Do đó $e=50-a-b-d\le 50-1-1-10=38$.
\step{3}{Ví dụ đạt được}
Danh sách $1,1,10,10,38$: tổng $60$, trung vị $10$. Vậy $e$ lớn nhất là $38$.

\needspace{6\baselineskip}\par\smallskip\textbf{Câu b): làm $e$ nhỏ nhất.}
\step{4}{Giới hạn dưới}
Muốn $e$ nhỏ thì $a,b$ phải lớn nhất có thể, nhưng $a\le b\le 10$, nên $a+b\le 20$. Suy ra $d+e=50-(a+b)\ge 30$. Vì $d\le e$ nên $2e\ge d+e\ge 30$, tức $e\ge 15$.
\step{5}{Ví dụ đạt được}
Danh sách $10,10,10,15,15$: tổng $60$, trung vị $10$. Vậy $e$ nhỏ nhất là $15$.

\needspace{6\baselineskip}\par\smallskip\textbf{Câu c): năm số đôi một khác nhau} ($a<b<10<d<e$).
\step{6}{$e$ lớn nhất}
Nhỏ nhất có thể: $a\ge 1$, $b\ge 2$, $d\ge 11$. Nên $e\le 50-1-2-11=36$. Ví dụ: $1,2,10,11,36$ (tổng $60$). Vậy $e$ lớn nhất là $36$.
\step{7}{$e$ nhỏ nhất}
Lớn nhất có thể: $b\le 9$, $a\le 8$, nên $a+b\le 17$ và $d+e\ge 50-17=33$. Vì $d<e$ nên $d\le e-1$, suy ra $2e-1\ge d+e\ge 33$, tức $e\ge 17$. Ví dụ: $8,9,10,16,17$ (tổng $60$). Vậy $e$ nhỏ nhất là $17$.

\begin{warnbox}Một bài ``lớn nhất/nhỏ nhất'' cần đủ \emph{hai} phần: (1) chứng minh không thể vượt qua giới hạn, và (2) một ví dụ cụ thể đạt đúng giới hạn. Chỉ có ví dụ thì chưa chắc đã là tốt nhất; chỉ có giới hạn thì chưa chắc đạt được.\end{warnbox}

\begin{parentbox}Có thể cho con dùng $60$ viên kẹo chia vào $5$ ``hộp'' xếp theo thứ tự, hộp giữa luôn có $10$ viên. Hỏi: muốn hộp cuối nhiều nhất thì các hộp khác nên có ít hay nhiều? Muốn hộp cuối ít nhất thì sao? Trò chơi này giúp con ``cảm'' được cách lập luận tối ưu.\end{parentbox}

\verify{$1+1+10+10+38=60$; $10+10+10+15+15=60$; $1+2+10+11+36=60$; $8+9+10+16+17=60$; tất cả đều có trung vị $10$. Máy tính duyệt mọi danh sách cũng xác nhận các giá trị $38$, $15$, $36$, $17$.}
\ans{a) $38$. \ b) $15$. \ c) Lớn nhất $36$, nhỏ nhất $17$.}
}
